%------------------------------------------------------------------------------%
\section{Revisão}

\begin{exercises*}

% Thomas pg 73 - ex 19-24
%------------------------------------------------------------------------------%
\begin{exercise}
  Calcule a soma das séries.
  \begin{mathlist}
    \item \sum_{n=3}^\infty \frac{1}{(2n\!-\!3)(2n\!-\!1)}
    \item \sum_{n=2}^\infty \frac{-2}{n(n+1)}
    \item \sum_{n=1}^\infty \frac{9}{(3n\!-\!1)(3n\!+\!2)}
    %
    \item \sum_{n=3}^\infty \frac{-8}{(4n\!-\!3)(4n\!+\!1)}
    \item \sum_{n=0}^\infty e^{-n}
    \item \sum_{n=1}^\infty \frac{3(-1)^n}{4^n}
  \end{mathlist}
\end{exercise}

% Thomas pg 73 - ex 19-24
%------------------------------------------------------------------------------%
\begin{exercise}
  Determine se as séries convergem absolutamente, convergem condicionalmente
  ou divergem.
  \begin{mathlist}
    % 25-32
    \item \sum_{n=1}^\infty \frac{1}{\sqrt{n}}
    \item \sum_{n=1}^\infty \frac{-5}{n}
    \item \sum_{n=1}^\infty \frac{(-1)^n}{\sqrt{n}}
    \item \sum_{n=1}^\infty \frac{1}{2n^3}
    \item \sum_{n=1}^\infty \frac{(-1)^n}{\ln(n+1)}
    \item \sum_{n=2}^\infty \frac{1}{n\big(\ln(n)\big)^2}
    \item \sum_{n=1}^\infty \frac{\ln(n)}{n^3}
    \item \sum_{n=3}^\infty \frac{\ln(n)}{\ln\!\big(\ln(n)\big)}
    % 33-40
    \item \sum_{n=1}^\infty \frac{(-1)^n}{n\sqrt{n^2+1}}
    \item \sum_{n=1}^\infty \frac{(-1)^n\,3\,n^2}{n^3+1}
    \item \sum_{n=1}^\infty \frac{n+1}{n!}
    \item \sum_{n=1}^\infty \frac{(-1)^n(n^2+1)}{2n^2+n-1}
    \item \sum_{n=1}^\infty \frac{(-3)^n}{n!}
    \item \sum_{n=1}^\infty \frac{2^n\,3^n}{n^n}
    \item \!\sum_{n=1}^\infty \!\!\frac{1}{\sqrt{n(n\!+\!1)(n\!+\!2)\!}}
    \item \sum_{n=2}^\infty \frac{1}{n\sqrt{n^2-1}}
  \end{mathlist}
\end{exercise}

% Exercício de prova do prof. André Pereira
%------------------------------------------------------------------------------%
\begin{exercise}
  Considere as sequências
  \[
    a_n = \ln\left(1+\frac{1}{n^2}\right)
    \qquad\text{e}\qquad
    b_n = \frac{1}{n^2}
  \]
  \vspace{-\baselineskip}
  \begin{alphalist}
    \item Quais das seguintes séries convergem e quais divergem?
          No item (i) use o teste da integral diretamente, não use que é uma série~$p$.
    \begin{romanlist}[2]
      \item \( \sum_{n=1}^\infty b_n \)
      \item \( \sum_{n=1}^\infty a_n \)
      \item \( \sum_{n=1}^\infty \frac{a_n}{b_n} \)
      \item \( \sum_{n=1}^\infty a_n^n \)
    \end{romanlist}
    \item Use o teste de Leibniz para mostrar que a série
    \( \sum_{n=1}^\infty(-1)^{n+1}a_n \)
    converge.
    \item Mostre que o teste da raiz não é conclusivo para a série
    \[
      \sum_{n=1}^\infty \left(\frac{a_n}{b_n}\right)^n =
      \sum_{n=1}^\infty \left[n^2\ln\left(1+\frac{1}{n^2}\right)\right]^n
    \]
    Por que o teste da raiz é inconclusivo neste caso? \\
    \hint{Dê exemplos de duas séries que satisfazem $\sqrt[n]{a_n}\rightarrow 1$,
      onde $a_n$ é o termo geral, mas uma converge e a outra diverge.}
  \end{alphalist}
\begin{answer}
  \begin{alphalist}
    \item
    \begin{minipage}[t]{0.9\linewidth}
      \begin{romanlist}
        \item \(\sum_{n=1}^\infty b_n=\sum_{n=1}^\infty \frac{1}{n^2}\)
        \item \(\sum_{n=1}^\infty a_n=\sum_{n=1}^\infty \ln\left(1+\frac{1}{n^2}\right)\)
        \item \(\sum_{n=1}^\infty \frac{a_n}{b_n}=\sum_{n=1}^\infty n^2\ln\left(1+\frac{1}{n^2}\right)\)
        \item \(\sum_{n=1}^\infty a_n^n=\sum_{n=1}^\infty \left[\ln\left(1+\frac{1}{n^2}\right)\right]^n\)
      \end{romanlist}
    \end{minipage}
%    \item ~
%    \item ~
  \end{alphalist}
\end{answer}
\end{exercise}

% Exercício de prova do prof. André Pereira
%------------------------------------------------------------------------------%
\begin{exercise}
  Considere as sequências
  \[
    a_n = \frac{n!}{5^n}    \hfill
    b_n = \frac{1}{n\ln n}  \hfill
    c_n = \left(\frac{1}{n}\right)^{\sfrac{1}{\ln(n)}} \hfill
  \]
  Use o teste indicado entre parenteses para verificar se as séries convergem ou divergem.
  No(s) caso(s) que converge(m) justifique se é absolutamente ou condicionalmente convergente.
  \begin{alphalist}
    \item \(\sum_{n=1}^\infty a_n\)           \tabto{35mm}(teste do termo geral)
    \item \(\sum_{n=1}^\infty\frac{1}{a_n}\)  \tabto{35mm}(teste da razão)
    \item \(\sum_{n=2}^\infty b_n\)           \tabto{35mm}(teste da integral)
    \item \(\sum_{n=2}^\infty (-1)^{n+1}b_n\) \tabto{35mm}(teste de Leibniz)
    \item \(\sum_{n=1}^\infty c_n\)           \tabto{35mm}(teste do termo geral)
  \end{alphalist}
\end{exercise}

% Exercício de prova do prof. André Pereira
%------------------------------------------------------------------------------%
\begin{exercise}
  Mostre que
  \begin{align*}
    \sum_{n=1}^\infty & \frac{\ln\left(\dfrac{n+1}{n+2}\right)}{\;\ln(n+2)\ln(n+1)\;}          \\
                      & = \sum_{n=1}^\infty \left(\frac{1}{\ln(n+2)}-\frac{1}{\ln(n+1)}\right)
  \end{align*}
  Encontre o limite desta série.
\end{exercise}

% Exercício de prova do prof. André Pereira
%------------------------------------------------------------------------------%
\begin{exercise}
  Considere as sequências
  \[
    a_n = \frac{(n!)^2}{(2n)!} \hfill
    b_n = \frac{1}{\sqrt{n}}   \hfill
    c_n = \frac{1}{2\sqrt{n}+\sqrt[3]{n}} \hfill
  \]
  Use o teste indicado entre parenteses para verificar se as séries
  convergem ou divergem. No(s) caso(s) que converge(m) justifique se é
  absolutamente ou condicionalmente convergente.
  \begin{alphalist}
    \item \(\sum_{n=1}^\infty a_n\)                \tabto{25mm}(teste da razão)
    \item \(\sum_{n=1}^\infty b_n\)                \tabto{25mm}(teste da integral)
    \item \(\sum_{n=1}^\infty b_n^{\sfrac{2}{n}}\) \tabto{25mm}(teste do termo geral)
    \item \(\sum_{n=2}^\infty c_n\)                \tabto{25mm}(teste da comparação no limite)
    \item \(\sum_{n=2}^\infty (-1)^{n+1}b_n\)      \tabto{35mm}(teste de Leibniz)
  \end{alphalist}
\end{exercise}

% Exercício de prova do prof. André Pereira
%------------------------------------------------------------------------------%
\begin{exercise}
  Encontre o limite da seguinte série
  \[
    \sum_{n=1}^\infty\frac{6}{(2n-1)(2n+1)}
  \]
\end{exercise}

% Exercício de prova do prof. André Pereira
%------------------------------------------------------------------------------%
\begin{exercise}
  Encontre o limite da seguinte série
  \[
    \sum_{n=1}^\infty\frac{6}{(2n-1)(2n+1)}
  \]
\end{exercise}

% Thomas pg 36 55-62
%------------------------------------------------------------------------------%
\begin{exercise}
Determine quais séries são convergentes ou divergentes.
\begin{mathlist}[2][series=ex-03-10-A]
  \item \sum_{n=1}^\infty  \frac{2^n\,n!\,n!}{(2n)!}
  \item \sum_{n=1}^\infty  \frac{(n!)^n}{n^{n^2}}
  \item \sum_{n=1}^\infty  \frac{n^n}{2^{n^2}}
  \item \sum_{n=1}^\infty  \frac{(n!)^n}{(n^n)^2}
  \item \sum_{n=1}^\infty  \frac{n^n}{(2^n)^2}
\end{mathlist}
\begin{mathlist}[1][resume=ex-03-10-A]
  \item \sum_{n=1}^\infty  \frac{(3n)!}{n!(n+1)!(n+2)!}
  \item \sum_{n=1}^\infty  \frac{1\cdot 3\cdots (2n-1)}{4^n\,2^n\,n!}
  \item \sum_{n=1}^\infty  \frac{1\cdot 3\cdots (2n-1)}{\big(2\cdot 4\cdots(2n)\big)\left(3^n+1\right)}
\end{mathlist}
\end{exercise}

% Thomas pg 74 1-4
%------------------------------------------------------------------------------%
\begin{exercise}
  Determine quais séries são convergentes ou divergentes.
  \begin{mathlist}[2]
    \item \sum_{n=1}^\infty \frac{1}{(3n-2)^{n+\sfrac{1}{2}}}
    \item \sum_{n=1}^\infty \frac{\big(\arctan(n)\big)^2}{n^2+1}
    \item \sum_{n=1}^\infty (-1)^n \tanh(n)
    \item \sum_{n=2}^\infty \frac{\log_n(n!)}{n!}
  \end{mathlist}
\end{exercise}

% Thomas pg 74 5-8
%------------------------------------------------------------------------------%
\begin{exercise}
  Considere as somas das sequências dadas a seguir, quais delas convergem?
  \begin{mathlist}
    \item a_1 = 1 \qquad a_{n+1} = \frac{n(n+1)a_n}{(n+2)(n+3)} \) \\
      \hint{escreva os primeiros termos e generalize os cancelamentos}
    \item a_1 = 7 \qquad a_{n+1} = \frac{n a_n}{(n+1)(n-1)}
    \item a_1 = 1 \qquad a_{n+1} = \frac{1}{1+a_n}
    \item a_n = \begin{cases}
      \dfrac{1}{3^n}, &\; n \text{ ímpar} \\[3mm]
      \dfrac{n}{3^n}, &\; n \text{ par}
    \end{cases}
  \end{mathlist}
\end{exercise}

%------------------------------------------------------------------------------%
\end{exercises*}
%------------------------------------------------------------------------------%
